\section*{Hoofdstuk \ref{ch:g1:light-and-shadows} --- Licht en schaduw}

\begin{solution}{exo:g1:light-and-shadows:1}
Bronnen: de kaarsvlam, de glimworm, de zaklamp. Geen bronnen: de
Maan, de spiegel, de witte muur --- die sturen alleen \omterm{def:g1:light-and-shadows:source}{licht} terug dat
erop valt.
\end{solution}

\begin{solution}{exo:g1:light-and-shadows:2}
\omterm{def:g1:light-and-shadows:source}{Licht} ontbreekt. Er is geen bron in de kamer, dus er valt geen
\omterm{def:g1:light-and-shadows:source}{licht} op de dingen, en dingen waar geen \omterm{def:g1:light-and-shadows:source}{licht} op valt kun je niet
zien.
\end{solution}

\begin{solution}{exo:g1:light-and-shadows:3}
Een \omterm{def:g1:light-and-shadows:source}{lichtbron}, iets dat het \omterm{def:g1:light-and-shadows:source}{licht} tegenhoudt, en een muur (of scherm,
of de grond) --- in die volgorde op één rechte lijn: de bron schijnt,
de blokkeerder houdt tegen, de muur laat de donkere plek zien.
\end{solution}

\begin{solution}{exo:g1:light-and-shadows:4}
De \omterm{def:g1:light-and-shadows:shadow}{schaduw} ligt rechts van je --- tegenover de Zon. Nadat je je
helemaal hebt omgedraaid, staat de Zon rechts van je, dus ligt de
\omterm{def:g1:light-and-shadows:shadow}{schaduw} links: altijd aan de kant die van de Zon af ligt.
\end{solution}

\begin{solution}{exo:g1:light-and-shadows:5}
Het speelgoed \emph{naar de lamp toe} schuiven maakt de \omterm{def:g1:light-and-shadows:shadow}{schaduw}
reusachtig; het \emph{naar de muur toe} schuiven laat de \omterm{def:g1:light-and-shadows:shadow}{schaduw}
krimpen tot de grootte van het speelgoed.
\end{solution}

\begin{solution}{exo:g1:light-and-shadows:6}
Nee. Een \omterm{def:g1:light-and-shadows:shadow}{schaduw} is geen ding dat ergens van gemaakt is --- het is een
plek waar het \omterm{def:g1:light-and-shadows:source}{licht} ontbreekt omdat iets het \omterm{def:g1:light-and-shadows:source}{licht} tegenhoudt. Waar
niets is, valt niets op te pakken.
\end{solution}

\begin{solution}{exo:g1:light-and-shadows:7}
De Zon is een veel sterkere bron dan welke lamp ook: haar
\omterm{def:g1:light-and-shadows:source}{licht} is zo sterk dat het de ogen voorgoed kan beschadigen, en ze
niet alleen even verblindt.
\end{solution}

\begin{solution}{exo:g1:light-and-shadows:8}
Laat in de middag, wanneer de Zon laag hangt. De tekening laat rond
het middaguur een korte \omterm{def:g1:light-and-shadows:shadow}{schaduw} aan de voet van de boom zien en 's
middags een lange, uitgerekte \omterm{def:g1:light-and-shadows:shadow}{schaduw}, allebei van de Zon af wijzend.
\end{solution}

\begin{solution}{exo:g1:light-and-shadows:9}
Sam heeft gelijk. De \omterm{def:g1:light-and-shadows:shadow}{schaduw} bleef de hele tijd precies tegenover de
Zon liggen. Mia draaide zich om, dus veranderde de \omterm{def:g1:light-and-shadows:shadow}{schaduw} van plaats
\emph{ten opzichte van haar eigen voor- en achterkant} --- maar op de
grond is hij nooit bewogen.
\end{solution}
